% !PRINTED 216
\weberpage{216}

\begin{center}
\S.\,57.

\gesperrt{Die Octaëdergruppe.}
\end{center}

Bei der Octaëdergruppe kommt ein System von sechs conjugirten
vierzähligen Polen vor. Wir haben demnach eine Substitution
4\textsuperscript{ten} Grades, deren Periode
\[
1, \ \Theta, \ \Theta^2, \ \Theta^3, \ \Theta^4 = 1
\]
ein Paar dieser conjugirten vierzähligen Pole giebt. Nach \S.\,53, 1.
können wir in der Gruppe die Substitutionen annehmen:
\weq{1}{\Theta(x) = ix, \quad \psi(x) = \frac{1}{x},}
und es ist
\weq{2}{\psi^2 = 1, \quad \Theta\psi = \psi\Theta^3, \quad \Theta^2\psi = \psi\Theta^2, \quad \Theta^3\psi = \psi\Theta.}
Die zu dem Systeme der vierzähligen Pole gehörige Grundform
6\textsuperscript{ten} Grades muss bis auf einen constanten Factor
ungeändert bleiben durch die Substitutionen $\Theta$ und $\psi$, und
sie muss ausserdem den Factor $x_1x_2$ enthalten, also von einer der
beiden Formen sein
\[
x_1x_2(x_1^4 - x_2^4), \quad x_1x_2(x_1^4 + x_2^4).
\]
Es ist gleichgültig, welche der beiden Annahmen wir verfolgen, da
die eine durch die Substitution $ix_2$ für $x_2$, was dem Uebergange
zu einer transformirten Gruppe entspricht, in die andere übergeht.
Nehmen wir
\weq{3}{f(x_1, x_2) = x_1x_2(x_1^4 - x_2^4)}
als Grundform 6\textsuperscript{ten} Grades an, so sind die sechs
vierzähligen Pole
\weq{4}{0, \ \infty, \ 1, \ -1, \ i, \ -i.}
Es muss nun nach \S.\,53, 2. in $G$ eine Substitution $\chi$ geben,
so dass
\[
\chi(-1) = 0, \quad \chi(+1) = \infty
\]
ist, und wenn man den constanten Factor wie bei der Tetraëdergruppe
bestimmt [\S.\,56, (2)], so findet sich
\weq{5}{\chi(x) = i\,\frac{x+1}{x-1}.}
Hieraus erhalten wir
\[
\chi^2(x) = \frac{x+i}{x-i}, \quad \chi^3(x) = x,
\]
%CONT>>
